\chapter{Badania}
\label{ch:badania}

W~rozdziale opisano przebieg eksperymentu przeprowadzonego zgodnie
z~metodyką przedstawioną w~rozdziale~\ref{ch:rozwiniecie}. Kolejne etapy,
wykorzystywane w~nich narzędzia oraz ich wyniki zestawiono
w~tabeli~\ref{tab:etapy}. Uzyskane wyniki omówiono
w~rozdziale~\ref{ch:wyniki}.

\begin{table}[ht]
\fcmtcaption{Etapy eksperymentu. Źródło: opracowanie własne.}\label{tab:etapy}
\centering\footnotesize%
\begin{tabular}{p{0.04\textwidth} p{0.33\textwidth} p{0.2\textwidth} p{0.3\textwidth}}
\toprule
Lp. & Etap & Narzędzie & Wynik \\
\midrule
1 & Generowanie zbioru uczącego i~trening sieci & Unet-Astar & wagi modelu U-Net \\
\midrule
2 & Utworzenie płytek testowych na podstawie schematów & EasyEDA & projekty PCB \\
\midrule
3 & Trasowanie płytek autorouterem & EasyEDA & wariant \emph{EasyEDA} \\
\midrule
4 & Eksport płytek do formatu Specctra & EasyEDA & pliki DSN \\
\midrule
5 & Trasowanie na podstawie plików DSN & Unet-Astar & ścieżki wariantów
  \emph{A*} i~\emph{Unet-A*}, statystyki wyszukiwania \\
  \midrule
6 & Eksport ścieżek & eksporter SES & pliki SES \\
\midrule
7 & Sklonowanie płytek i~naniesienie wyników & EasyEDA & płytki wariantów
  \emph{A*} i~\emph{Unet-A*} \\
  \midrule
8 & Odczyt parametrów płytek & EasyEDA & miary jakości trasowania \\
\midrule
9 & Zestawienie wyników & arkusz / skrypt & tabele porównawcze \\
\bottomrule
\end{tabular}
\end{table}

\section{Przygotowanie sieci neuronowej}
\label{sec:exp-trening}

Zbiór uczący wygenerowano zgodnie z~opisem w~sekcji~\ref{sec:zbior-uczacy}.
Utworzono 1250~losowych płytek dwuwarstwowych o~bokach od 1000 do
2000~mil. Generowanie prowadzono równolegle w~pięciu procesach.
Z~płytek uzyskano łącznie 13830~próbek (średnio ok.~11 na płytkę),
zajmujących na dysku 3,1~GB. Zbiór podzielono losowo w~proporcji 9:1 na
część uczącą i~walidacyjną.

Sieć U-Net trenowano z~parametrami podanymi w~tabeli~\ref{tab:trening}.
Do dalszych badań wybrano model o~najmniejszej stracie na zbiorze
walidacyjnym. Trening na karcie graficznej trwał łącznie 13~h 27~min
(15~epok, średnio ok.~54~min na epokę). Strata walidacyjna malała niemal
w~każdej epoce, z~0,155 po pierwszej do 0,098 po ostatniej. Najmniejszą
stratę walidacyjną (0,0978, w~tym BCE 0,0486 i~strata Dice 0,1471)
osiągnięto w~ostatniej, 15.~epoce. Oznacza to, że dłuższy trening mógłby
jeszcze nieznacznie poprawić model.
Przebieg funkcji straty przedstawiono w~rozdziale~\ref{ch:wyniki}.

\section{Przygotowanie płytek testowych}
\label{sec:exp-plytki}

Płytki testowe utworzono w~EasyEDA na podstawie dwóch rodzajów schematów:
\begin{itemize}
  \item schematów własnych,
  \item schematu publicznego projektu przykładowego \emph{Project-Sample},
    udostępnionego w~sekcji \emph{Modules} serwisu
    EasyEDA~\cite{EasyEDA_ProjectSample}.
\end{itemize}

Dla każdego schematu wygenerowano w~EasyEDA płytkę PCB, na której
umieszczono footprinty elementów wraz z~siecią połączeń. Następnie
wyznaczono prostokątny obrys płytki i~rozmieszczono elementy
ręcznie. Rozmieszczenie elementów nie było zmieniane w~dalszych etapach, 
dlatego wszystkie warianty trasowania otrzymują identyczne położenia padów. 
Płytki spełniają wymagania podane w~sekcji~\ref{sec:zbior-testowy}:
mają dwie warstwy miedzi, prostokątny obrys i~nie zawierają 
stref miedzi ani obszarów zabronionych.

Dla wszystkich płytek i~wszystkich sieci zastosowano ten sam zestaw reguł
projektowych (tabela~\ref{tab:reguly}). Reguły te zapisywane są w~pliku DSN
i~obowiązują wszystkie trzy warianty trasowania. W~programie Unet-Astar są
one zaokrąglane w~górę do całkowitej liczby oczek siatki
(sekcja~\ref{sec:format-danych}).

\begin{table}[ht]
\fcmtcaption{Reguły projektowe płytek testowych. Źródło: opracowanie własne.}\label{tab:reguly}
\centering\footnotesize%
\begin{tabular}{l r r}
\toprule
Reguła (wartość minimalna) & [mm] & [mil] \\
\midrule
Szerokość ścieżki & 0{,}254 & 10 \\
\midrule
Odstęp izolacyjny (\emph{clearance}) & 0{,}152 & 6 \\
\midrule
Średnica przelotki & 0{,}610 & 24 \\
\midrule
Średnica otworu przelotki & 0{,}305 & 12 \\
\bottomrule
\end{tabular}
\end{table}

Charakterystykę płytek testowych zestawiono w~tabeli~\ref{tab:plytki-testowe}.

\begin{table}[ht]
\fcmtcaption{Płytki testowe. Źródło: opracowanie własne.}\label{tab:plytki-testowe}
\centering\footnotesize%
\begin{tabular}{l l r r r r}
\toprule
Płytka & Pochodzenie & Wymiary [mm] & Elementy & Sieci & Połączenia \\
\midrule
Płytka z~diodą LED & własna & $51 \times 34$ & 4 & 4 & 4 \\
\midrule
Płytka przykładowa & publiczna & $37 \times 25$ & 8 & 6 & 16 \\
\bottomrule
\end{tabular}
\end{table}

Kolumna \emph{Połączenia} podaje liczbę połączeń dwupinowych, na które
program Unet-Astar dzieli sieci wielopinowe. Odpowiada ona liczbie
zadań wyszukiwania ścieżki wykonywanych przez algorytm A*. Na płytce
przykładowej jedna z~sześciu sieci ma tylko jeden pin, więc nie wymaga
poprowadzenia. Pozostałe pięć sieci, w~tym ośmiopinowa sieć zasilania
+24~V, daje łącznie 16~połączeń.

\section{Trasowanie autorouterem EasyEDA}
\label{sec:exp-autorouter}

Każdą płytkę testową poprowadzono autorouterem wbudowanym w~EasyEDA
(EasyEDA Standard 6.5.61, tabela~\ref{tab:srodowisko}). Reguły trasowania
autoroutera (\emph{Routing Rules}) odwołują się do reguł projektowych
(\emph{Design Rules}) zdefiniowanych dla projektu. Autorouter korzystał
więc z~tych samych wartości szerokości ścieżki, odstępu izolacyjnego
i~wymiarów przelotki, które podano w~tabeli~\ref{tab:reguly}.
Poprowadzona płytka stanowi wariant odniesienia \emph{EasyEDA}.
Dla każdej płytki zanotowano, czy autorouter poprowadził wszystkie
połączenia. EasyEDA nie podaje czasu działania autoroutera, dlatego czasu
trasowania tego wariantu nie mierzono.

\section{Eksport plików DSN}
\label{sec:exp-dsn}

Płytki wyeksportowano z~EasyEDA do plików w~formacie \emph{Specctra DSN}
(sekcja~\ref{sec:format-danych}). Plik zawiera obrys płytki, warstwy,
reguły projektowe, kształty padów, położenia pinów i~listę sieci. Zawiera
również ścieżki poprowadzone przez autorouter (sekcja \emph{wiring}). Program
Unet-Astar pomija tę sekcję, więc warianty \emph{A*} i~\emph{Unet-A*}
trasują płytkę od początku, a~wynik autoroutera nie wpływa na ich działanie.

\section{Trasowanie programem Unet-Astar}
\label{sec:exp-unet-astar}

Na podstawie każdego pliku DSN płytkę poprowadzono w~dwóch wariantach:
\begin{description}
  \item[A*] --- algorytm wielowarstwowego A* bez udziału sieci,
  \item[Unet-A*] --- ten sam algorytm wspomagany predykcją wytrenowanej
    sieci, z~parametrami podanymi w~sekcji~\ref{sec:trening}.
\end{description}
W~wariancie \emph{Unet-A*} każde połączenie korzystające z~sieci było
dodatkowo wyznaczane bez jej udziału, przy tym samym stanie płytki. Pozwala
to porównać liczbę rozwiniętych węzłów i~czas wyszukiwania dla tych samych
zadań, niezależnie od kolejności i~przebiegu wcześniejszych połączeń.
Dla każdej płytki zapisano liczbę rozwiniętych węzłów, czas trasowania,
długość ścieżek, liczbę przelotek i~liczbę poprowadzonych połączeń.

Czas trasowania zmierzono osobno dla obu wariantów, w~identyczny sposób.
Jest to czas działania procedury trasowania: od podziału sieci na
połączenia dwupinowe do wyznaczenia ostatniej ścieżki. Nie obejmuje on
wczytania pliku DSN, załadowania modelu ani zapisu pliku SES.
W~wariancie \emph{Unet-A*} czas obejmuje inferencję sieci, a~kontrolne
wyszukiwanie bez sieci opisane wyżej było przy tym pomiarze wyłączone.
Każdy wariant uruchomiono trzykrotnie i~podano medianę. Przed pomiarami
wykonano jedno nieliczone uruchomienie, aby wyeliminować wpływ
inicjalizacji karty graficznej.

Dla płytek mieszczących się w~wejściu sieci wygenerowano ponadto próbki
w~sposób identyczny jak dla zbioru uczącego. Posłużyły one do wyznaczenia
miar Dice i~IoU (tabela~\ref{tab:miary}), opisujących jakość predykcji sieci
na płytkach rzeczywistych.

\section{Eksport ścieżek}
\label{sec:exp-ses}

Ścieżki wyznaczone w~wariantach \emph{A*} i~\emph{Unet-A*} zapisano
eksporterem opisanym w~sekcji~\ref{sec:format-danych}. Dla każdej płytki
powstały dwa pliki \emph{Specctra SES}, po jednym dla każdego wariantu.
Pliki odwołują się do pliku DSN, z~którego powstały, i~zawierają wyłącznie
ścieżki i~przelotki --- położenia elementów pozostają bez zmian.

\section{Naniesienie wyników na płytki w~EasyEDA}
\label{sec:exp-import}

W~EasyEDA każdą płytkę testową sklonowano dwukrotnie, osobno dla wariantu
\emph{A*} i~\emph{Unet-A*}. Ze sklonowanych płytek usunięto ścieżki
i~przelotki poprowadzone przez autorouter, pozostawiając elementy w~tych
samych położeniach. Następnie do każdej kopii zaimportowano odpowiedni plik
SES. W~ten sposób dla każdej płytki testowej uzyskano trzy projekty
różniące się wyłącznie poprowadzeniem ścieżek. Zaimportowane płytki
porównano wizualnie z~wynikiem programu Unet-Astar, aby wykluczyć błędy
przeliczenia współrzędnych.

\section{Odczyt parametrów płytek}
\label{sec:exp-odczyt}

Dla każdego z~trzech projektów płytki odczytano w~EasyEDA:
\begin{itemize}
  \item łączną długość ścieżek,
  \item liczbę przelotek,
  \item liczbę niepoprowadzonych połączeń,
  \item liczbę i~rodzaje naruszeń reguł projektowych zgłoszonych przez
    DRC.
\end{itemize}
% TODO: nazwa funkcji / okna EasyEDA, z którego odczytywane są parametry.
Parametry jakości trasowania odczytano tym samym narzędziem dla wszystkich
wariantów. Dzięki temu nie zależą one od sposobu liczenia długości ścieżek
w~programie Unet-Astar, który posługuje się siatką o~boku 8~mil.
Miary wydajności wyszukiwania (liczba rozwiniętych węzłów, czas) pochodzą
z~programu Unet-Astar i~dotyczą wyłącznie wariantów \emph{A*}
i~\emph{Unet-A*}.

\section{Zestawienie wyników}
\label{sec:exp-zestawienie}

Parametry odczytane dla wszystkich płytek i~wariantów zestawiono w~tabelach
porównawczych. Dla każdej płytki porównano warianty \emph{EasyEDA},
\emph{A*} i~\emph{Unet-A*} pod względem miar jakości trasowania, a~warianty
\emph{A*} i~\emph{Unet-A*} --- dodatkowo pod względem liczby rozwiniętych
węzłów i~czasu wyszukiwania. Wyniki zagregowano również dla całego zbioru
testowego. Wyniki przedstawiono i~omówiono w~rozdziale~\ref{ch:wyniki}.
